% !TEX root = ../main.tex
% =========================================================
% bab1.tex - PENDAHULUAN
% =========================================================

\chapter[PENDAHULUAN]{\\ PENDAHULUAN}

\section{Latar Belakang}

Pemanfaatan energi terbarukan menjadi salah satu aspek penting dalam
pengembangan sistem ketenagalistrikan modern. Salah satu teknologi yang
berkembang pesat adalah pembangkit listrik tenaga surya (PLTS). Selain dibangun
pada lahan darat, PLTS juga dapat dikembangkan pada permukaan air dalam bentuk
PLTS terapung. PLTS terapung memiliki keuntungan dari sisi pemanfaatan area
perairan, tetapi memiliki tantangan operasi dan pemeliharaan yang berbeda
dibandingkan PLTS berbasis darat karena komponen sistem berada pada lingkungan
yang dipengaruhi air, kelembapan, perubahan muka air, serta pergerakan struktur
apung \cite{worldbankFPV2019,ieaFPV2025}. PLTS Terapung Cirata sebagai lokasi
kerja praktik ditunjukkan pada Gambar~\ref{gbr:plts-terapung-cirata}.

\begin{figure}[H]
\centering
\includegraphics[width=0.82\textwidth]{pltscirata.jpg}
\caption{PLTS Terapung Cirata.}
\label{gbr:plts-terapung-cirata}
\end{figure}

Salah satu kondisi abnormal yang dapat ditemukan pada PLTS terapung adalah
modul PV yang terendam akibat gangguan pada struktur apung. Kondisi tersebut
merupakan gejala fisik yang perlu dievaluasi melalui parameter elektrik, yaitu
arus, tegangan, serta karakteristik kurva arus--tegangan (I--V) dan
daya--tegangan (P--V). Studi pada \textit{photovoltaic power plant} yang
mengalami kondisi banjir menunjukkan bahwa paparan air dan kelembapan dapat
berkaitan dengan penurunan tahanan isolasi serta munculnya arus bocor pada
modul PV \cite{ketjoy2022}. Oleh karena itu, modul yang pernah terendam tidak
dapat langsung dinyatakan layak hanya berdasarkan inspeksi visual. Alur
evaluasi modul PV pasca-terendam pada laporan ini ditunjukkan pada
Gambar~\ref{gbr:konteks-latarbelakang}.

\begin{figure}[H]
\centering
\begin{tikzpicture}[
    node distance=0.75cm,
    every node/.style={font=\small},
    box/.style={
        draw,
        rounded corners=3pt,
        align=center,
        minimum width=8.6cm,
        minimum height=0.78cm,
        fill=blue!6
    },
    boxwarn/.style={
        draw,
        rounded corners=3pt,
        align=center,
        minimum width=8.6cm,
        minimum height=0.78cm,
        fill=orange!12
    },
    boxresult/.style={
        draw,
        rounded corners=3pt,
        align=center,
        minimum width=8.6cm,
        minimum height=0.78cm,
        fill=green!10
    },
    arrow/.style={
        -{Latex[length=2.3mm]},
        line width=0.7pt
    }
]

\node[box] (a) {PLTS Terapung};
\node[box, below=of a] (b) {Lingkungan air, kelembapan, perubahan muka air, dan pergerakan
struktur apung};
\node[boxwarn, below=of b] (c) {Modul PV pasca-terendam};
\node[box, below=of c] (d) {Pengujian kurva I--V dan koreksi ke STC};
\node[box, below=of d] (e) {Analisis $P_{max}$, Fill Factor, $R_{VM}$, dan $R_{IM}$};
\node[boxresult, below=of e] (f) {Evaluasi performa dan arah pemeriksaan O\&M};

\draw[arrow] (a) -- (b);
\draw[arrow] (b) -- (c);
\draw[arrow] (c) -- (d);
\draw[arrow] (d) -- (e);
\draw[arrow] (e) -- (f);

\end{tikzpicture}
\caption{Alur evaluasi modul PV pasca-terendam.}
\label{gbr:konteks-latarbelakang}
\end{figure}

Pengujian kelayakan modul PV pada umumnya dilakukan melalui pengukuran kurva
I--V. Karakteristik I--V mengandung informasi mengenai keadaan aktual modul,
termasuk parameter \textit{open-circuit voltage} ($V_{oc}$),
\textit{short-circuit current} ($I_{sc}$), titik daya maksimum ($P_{max}$),
serta titik kerja $V_{mp}$ dan $I_{mp}$ \cite{linwu2025}. Parameter tersebut
dibandingkan dengan nilai \textit{nameplate} pada \textit{datasheet} untuk
mengevaluasi tingkat penurunan performa modul. Perubahan bentuk kurva I--V juga
dipengaruhi oleh perubahan parameter internal sel surya, seperti resistansi
seri ($R_s$) dan resistansi shunt ($R_{sh}$), sehingga bentuk kurva dapat
digunakan sebagai petunjuk awal kecenderungan degradasi modul
\cite{villalva2009}. Keterkaitan antara fenomena fisik akibat paparan air dan
parameter elektrik yang diamati ditunjukkan pada
Gambar~\ref{gbr:hubungan-fisik-elektrik}.

\begin{figure}[H]
\centering
\begin{tikzpicture}[
    node distance=0.9cm,
    every node/.style={font=\small},
    mainbox/.style={
        draw,
        rounded corners=3pt,
        align=center,
        minimum width=3.7cm,
        minimum height=0.85cm,
        fill=blue!6
    },
    warnbox/.style={
        draw,
        rounded corners=3pt,
        align=center,
        minimum width=3.7cm,
        minimum height=0.85cm,
        fill=orange!12
    },
    resultbox/.style={
        draw,
        rounded corners=3pt,
        align=center,
        minimum width=3.7cm,
        minimum height=0.85cm,
        fill=green!10
    },
    arrow/.style={
        -{Latex[length=2.3mm]},
        line width=0.7pt
    }
]

\node[warnbox] (a) at (0,0) {Paparan air dan\\ kelembapan};
\node[mainbox, right=1.05cm of a] (b) {Perubahan kontak,\\ isolasi, atau\\ resistansi internal};
\node[mainbox, right=1.05cm of b] (c) {Perubahan $V_{oc}$,\\ $I_{sc}$, $V_{mp}$,\\ dan $I_{mp}$};

\node[resultbox, below=1.0cm of b] (d) {Perubahan $P_{max}$,\\ Fill Factor,\\ $R_{VM}$, dan $R_{IM}$};

\node[resultbox, below=1.0cm of d] (e) {Indikasi awal\\ degradasi modul};

\draw[arrow] (a) -- (b);
\draw[arrow] (b) -- (c);
\draw[arrow] (b) -- (d);
\draw[arrow] (c) |- (d);
\draw[arrow] (d) -- (e);

\end{tikzpicture}
\caption{Hubungan fenomena fisik dan parameter elektrik modul PV.}
\label{gbr:hubungan-fisik-elektrik}
\end{figure}

Untuk memisahkan informasi sisi tegangan dan sisi arus dari kurva I--V, Pei dan
Hao \cite{pei2019} mendefinisikan indikator $R_{VM} = V_{mp}/V_{oc}$ dan
$R_{IM} = I_{mp}/I_{sc}$. Indikator $R_{VM}$ mencerminkan posisi relatif titik
daya maksimum terhadap $V_{oc}$, sedangkan $R_{IM}$ mencerminkan posisi relatif
titik daya maksimum terhadap $I_{sc}$. Keduanya dapat didekomposisi dari Fill
Factor melalui hubungan $FF = R_{VM} \cdot R_{IM}$, sehingga penurunan $FF$
dapat diidentifikasi lebih lanjut apakah dominan berasal dari sisi tegangan
atau sisi arus. Pendekatan serupa telah diterapkan pada PLTS skala 1~MWp untuk
membedakan pola perubahan parameter arus dan tegangan pada level sistem
\cite{sugirianta2022}.

Modul PV pada sistem PLTS tidak bekerja sebagai komponen tunggal. Beberapa
modul tersusun seri membentuk \textit{string} dan beberapa \textit{string}
dihubungkan paralel menuju \textit{DC combiner} dan inverter. Kondisi abnormal
pada satu modul dapat memengaruhi karakteristik elektrik \textit{string},
baik melalui penurunan kontribusi arus maupun penurunan kontribusi tegangan.
Indikator $R_{VM}$ dan $R_{IM}$ membantu membedakan kedua kecenderungan pola
tersebut tanpa harus menetapkan satu jenis gangguan secara final.

Selain aspek performa daya, kondisi modul PV terendam juga berkaitan dengan
keselamatan dan isolasi. Parameter arus dan tegangan dari pengujian I--V belum
cukup untuk menyimpulkan bahwa suatu modul berada pada kondisi aman dari sisi
isolasi. Prosedur pemeriksaan \textit{ground fault} menempatkan pengukuran
tahanan isolasi sebagai langkah penting ketika pengukuran tegangan belum cukup
memberikan informasi yang jelas \cite{smaGroundFault}. Oleh sebab itu, hasil
pengujian I--V diposisikan sebagai dasar elektrik awal, sedangkan keputusan
pemasangan kembali tetap memerlukan \textit{insulation resistance test} dan
inspeksi lanjutan.

Topik ini dipilih karena relevan dengan kegiatan \textit{operation and
maintenance} (O\&M) PLTS Terapung Cirata. Selama pelaksanaan kerja praktik,
sejumlah modul PV yang mengalami kondisi abnormal telah dibawa ke darat dan
diuji menggunakan \textit{I--V curve tracer} dengan koreksi parameter ke
kondisi \textit{standard test conditions} (STC). Koreksi ke STC diperlukan
agar parameter antar modul dapat dibandingkan pada kondisi acuan yang sama,
dengan tetap memperhatikan batasan ketidakpastian pengukuran pada irradiansi
tertentu \cite{iec60891,iec61829}. Analisis pada laporan ini dilakukan terhadap
59 data pengujian modul PV tipe JKM560N-72HL4-BDV dan diposisikan sebagai
evaluasi awal berbasis parameter elektrik yang memanfaatkan perbandingan
parameter STC terhadap \textit{nameplate}, Fill Factor, indikator $R_{VM}$ dan
$R_{IM}$, serta interpretasi awal terhadap implikasi O\&M.

\section{Batasan Masalah}

Agar pembahasan dalam laporan kerja praktik ini lebih terarah, batasan masalah
yang digunakan adalah sebagai berikut:

\begin{enumerate}
    \item Analisis dilakukan pada konteks PLTS Terapung Cirata.
    \item Objek analisis adalah 59 data pengujian kurva I--V modul PV tipe
          JKM560N-72HL4-BDV yang sebelumnya mengalami kondisi abnormal
          (pasca-terendam).
    \item Modul PV terendam dipandang sebagai kondisi abnormal fisik yang perlu
          dievaluasi melalui parameter elektrik; analisis tidak menetapkan
          satu jenis gangguan secara final.
    \item Analisis menggunakan parameter STC, deviasi terhadap
          \textit{nameplate}, Fill Factor, serta indikator $R_{VM}$ dan
          $R_{IM}$.
    \item Pengujian I--V dilakukan secara \textit{onshore} dengan koreksi
          parameter ke STC; analisis ini tidak mencakup pelaksanaan pengujian
          tahanan isolasi.
    \item Analisis tidak membahas kestabilan sistem tenaga, profil frekuensi,
          maupun profil tegangan pada jaringan AC.
    \item Hasil analisis merupakan evaluasi awal berbasis parameter elektrik
          dan bukan pengganti prosedur resmi operasi dan pemeliharaan
          perusahaan.
\end{enumerate}

\section{Tujuan}

Berikut adalah beberapa tujuan yang ditetapkan untuk memandu penyusunan laporan
kerja praktik ini:

\begin{enumerate}
    \item Menjelaskan relevansi kondisi modul PV pasca-terendam terhadap
          kebutuhan evaluasi berbasis parameter elektrik pada PLTS terapung.
    \item Menganalisis 59 data pengujian I--V melalui perbandingan parameter
          STC terhadap nilai \textit{nameplate}.
    \item Menghitung dan menganalisis Fill Factor serta indikator $R_{VM}$ dan
          $R_{IM}$ untuk mengidentifikasi pola degradasi yang dominan.
    \item Mengidentifikasi kecenderungan pola degradasi berdasarkan indikator
          arus--tegangan, yaitu apakah penurunan performa lebih dominan pada
          sisi tegangan atau sisi arus.
    \item Menyusun interpretasi awal sebagai dasar arah pemeriksaan teknis
          pada kegiatan \textit{operation and maintenance}.
\end{enumerate}

\section{Manfaat}

Adapun manfaat yang diperoleh dari penyusunan laporan kerja praktik ini adalah
sebagai berikut:

\begin{enumerate}
    \item Bagi penulis, kegiatan kerja praktik ini memberikan pemahaman
          mengenai konfigurasi modul, \textit{string}, dan \textit{array} pada
          PLTS terapung serta penerapan pengujian I--V dan indikator arus--%
          tegangan untuk mengevaluasi modul PV pasca-terendam.
    \item Bagi perusahaan, laporan ini dapat memberikan kerangka awal untuk
          memilah modul PV pasca-terendam berdasarkan performa daya dan
          mengarahkan pemeriksaan lanjutan. Jika indikator tegangan lebih
          dominan menurun, pemeriksaan dapat difokuskan pada sisi tegangan
          operasi; jika hasil I--V belum cukup menjelaskan kondisi modul,
          pemeriksaan dapat dilanjutkan melalui \textit{insulation resistance
          test} dan inspeksi visual.
    \item Bagi akademik, laporan ini dapat menjadi contoh penerapan teori
          karakteristik I--V/P--V serta analisis parameter elektrik modul PV
          pada kasus PLTS terapung.
\end{enumerate}

\section{Rumusan Masalah}

Berdasarkan latar belakang yang telah dijelaskan, rumusan masalah dalam laporan
kerja praktik ini adalah sebagai berikut:

\begin{enumerate}
    \item Bagaimana status performa daya 59 data pengujian modul PV
          pasca-terendam berdasarkan $P_{max}$ STC?
    \item Bagaimana deviasi parameter $V_{oc}$, $I_{sc}$, $V_{mp}$, $I_{mp}$,
          dan $P_{max}$ terhadap nilai \textit{nameplate}?
    \item Bagaimana Fill Factor serta indikator $R_{VM}$ dan $R_{IM}$
          membedakan kelompok Ok dan Not Ok, dan pola degradasi mana yang lebih
          dominan?
    \item Bagaimana hasil analisis tersebut dapat digunakan sebagai dasar
          interpretasi awal untuk kegiatan \textit{operation and maintenance}?
\end{enumerate}

\section{Waktu dan Tempat Pelaksanaan}

Pelaksanaan kerja praktik dilakukan pada periode yang ditetapkan oleh
perusahaan dan program studi. Adapun waktu dan tempat pelaksanaan kerja praktik
ditunjukkan pada Tabel~\ref{tabel:waktu-tempat-kp}.

\begin{table}[h!]
    \centering
    \begin{tabular}{p{1cm} p{2.5cm} p{0.25cm} p{10.25cm}}
        & Waktu  & : & {\tglMulai} -- {\tglSelesai} \\
        & Tempat & : & {\penempatan} \\
    \end{tabular}
    \caption{Waktu dan tempat pelaksanaan kerja praktik}
    \label{tabel:waktu-tempat-kp}
\end{table}

\section{Metode}

Metode yang digunakan dalam penyusunan laporan kerja praktik ini meliputi
observasi kegiatan kerja praktik, studi literatur, pengumpulan data pengujian
yang diperbolehkan, pengolahan data, serta analisis dan interpretasi hasil.
Observasi dilakukan untuk memahami lingkungan kerja serta kegiatan operasi dan
pemeliharaan di PLTS Terapung Cirata. Studi literatur dilakukan untuk memahami
konfigurasi sistem PV, karakteristik I--V/P--V, parameter STC, Fill Factor,
indikator arus--tegangan, serta peran \textit{insulation resistance test} pada
modul pasca-terendam. Data yang dianalisis berupa 59 data pengujian kurva I--V
modul PV pasca-terendam yang telah dikoreksi ke STC. Pengolahan data dilakukan
dalam empat tahap, yaitu perbandingan parameter terhadap \textit{nameplate},
analisis Fill Factor, analisis indikator $R_{VM}$ dan $R_{IM}$, serta
interpretasi pola degradasi dan implikasinya terhadap kegiatan O\&M.

\section{Lingkup Penulisan}

Laporan kerja praktik ini disusun dalam lima bab. BAB~I berisi pendahuluan yang
meliputi latar belakang, batasan masalah, tujuan, manfaat, rumusan masalah,
waktu dan tempat pelaksanaan, metode, serta lingkup penulisan. BAB~II berisi
profil PT Pembangkitan Jawa Bali Masdar Solar Energi (PT~PMSE) dan gambaran
umum PLTS Terapung Cirata. BAB~III berisi pokok pembahasan yang mencakup dasar
teori dan metodologi teknis, meliputi sistem PLTS terapung, konfigurasi
modul-\textit{string}-\textit{array}, karakteristik I--V/P--V, parameter
elektrikal, Fill Factor, indikator arus--tegangan, pola perubahan parameter
arus dan tegangan pada sistem PV, serta objek dan tahapan pengolahan data.
BAB~IV berisi hasil dan pembahasan yang mencakup karakteristik data pengujian,
analisis status performa, deviasi parameter, Fill Factor, indikator $R_{VM}$
dan $R_{IM}$, interpretasi pola degradasi, serta implikasi terhadap kegiatan
\textit{operation and maintenance}. BAB~V berisi kesimpulan dan saran.